\documentclass[11pt,a4paper]{article}

\usepackage[utf8]{inputenc}
\usepackage[T1]{fontenc}
\usepackage[spanish,es-tabla,es-noquoting]{babel}
\usepackage{lmodern}
\usepackage{microtype}
\usepackage[a4paper,top=2.5cm,bottom=2.5cm,left=2.5cm,right=2.2cm]{geometry}
\usepackage[table]{xcolor}
\usepackage{tabularx}
\usepackage{array}
\usepackage{ragged2e}
\usepackage{amsmath}
\usepackage{xurl}
\usepackage[hidelinks]{hyperref}

\definecolor{uteqBlue}{HTML}{0E3F7E}
\definecolor{uteqGray}{HTML}{F3F4F6}

\setlength{\parindent}{0pt}
\setlength{\parskip}{5pt}
\renewcommand{\arraystretch}{1.3}

\newcolumntype{L}[1]{>{\RaggedRight\arraybackslash}p{#1}}
\newcolumntype{C}[1]{>{\centering\arraybackslash}p{#1}}
\newcolumntype{Y}{>{\RaggedRight\arraybackslash}X}
\newcommand{\hdr}[1]{\textcolor{white}{\textbf{#1}}}

\begin{document}

\begin{center}
    {\Large\bfseries\color{uteqBlue} Informe de coevaluación de exposiciones}\\[4pt]
    {\large Aplicaciones Distribuidas (ISR-701) --- GA-EXP-02 - FORM}\\[4pt]
    Evaluador: Zambrano Parraga Luis Mario\\
    Equipo: MPSZ (tema T5)\\
    Fecha de la sesión: 26/11/2026
\end{center}

%============================================================
% GRUPO 1
%============================================================

\section{Equipo GGHL --- Computación paralela}

\textbf{Hora de inicio:} [08:30] \quad
\textbf{Hora de fin:} [09:00] \quad
\textbf{Duración:} [30] min \quad
\textbf{Tiempo asignado $T$:} [30] min

\begin{tabularx}{\linewidth}{|C{0.9cm}|L{3.4cm}|C{1.1cm}|Y|}
    \hline
    \rowcolor{uteqBlue}
    \hdr{D} & \hdr{Dimensión (peso)} & \hdr{Nivel} &
    \hdr{Evidencia concreta observada} \\
    \hline

    D1 & Cumplimiento de la guía (10) & [2--4] &
    Faltó hilos, sincronización y demás. \\
    \hline

    D2 & Fundamentación conceptual (20) & [2--4] &
    Se omitio el concepto de hilos, porque falto uno de los integrantes\\
    \hline

    D3 & Desarrollo técnico y demostración (15) & [2--4] &
    Diapositiva 8: Roger solo explicó una imagen de código,
    mas no vi ejecución local. Ejercicio de Bryan:
    explica el código y lógica, pero no introdujo un poco
    a qué PFC se está basando. \\
    \hline

    D4 & Análisis crítico y recomendación (20) & [2--4] &
    Estuvo limitado porque falto la explicacion sobre hilos, y de recomendación a roger que explique en consola, no solo en imagen para ver ese proceso en la practica\\
    \hline

    D5 & Comunicación oral y visual (25) & [3--4] &
    Demoraron un poco de tiempo preparando la exposición.\\
    \hline

    D6 & Fuentes y créditos (10) & [2--4] &
    Diapositiva 1: faltan fuentes que están como pendientes
    en la presentación. \\
    \hline

\end{tabularx}

\textbf{Nota sugerida:}

\[
    \frac{
        10(2)+20(2)+15(2)+20(2)+25(3)+10(2)
    }{4}
    \times \frac{10}{100}
    = 5.625 \approx 5.6
\]

\textbf{Fortaleza:} Bryan explico muy bien el código y la logica, primero en imagen y luego aplicado en codigo.

\textbf{Mejora:} Podrian mejorar la explicacion de hilos y sincronizacion, ya que no se explico en la exposicion, explicar el funcionamiento o recorrido del codigo en vez de presentar una sola imagen.

\textbf{Hoja de campo:}
\path{hojas_campos/image.png}

\clearpage

%============================================================
% GRUPO 2 - EVALUACION PENDIENTE DE FINALIZAR
%============================================================

\section{Equipo LMRT --- RPC, gRPC, Java RMI y modelos de comunicación aplicados a propuestas de distribución de los proyectos del equipo.}

\textbf{Hora de inicio:} [09:00] \quad
\textbf{Hora de fin:} [09:30] \quad
\textbf{Duración:} [30] min \quad
\textbf{Tiempo asignado $T$:} [30] min

\begin{tabularx}{\linewidth}{|C{0.9cm}|L{3.4cm}|C{1.1cm}|Y|}
    \hline
    \rowcolor{uteqBlue}
    \hdr{D} & \hdr{Dimensión (peso)} & \hdr{Nivel} &
    \hdr{Evidencia concreta observada} \\
    \hline

    D1 & Cumplimiento de la guía (10) & [0--4] &
    Todavia falta una de exponer un integrante\\
    \hline

    D2 & Fundamentación conceptual (20) & [0--4] &
    Todavia falta una de exponer un integrante\\
    \hline

    D3 & Desarrollo técnico y demostración  (15) & [0--4] &
    Todavia falta una de exponer un integrante \\
    \hline

    D4 & Análisis crítico y recomendación (20) & [0--4] &
    Todavia falta una de exponer un integrante\\
    \hline

    D5 & Comunicación oral y visual (25) & [2--4] &
    Demora en general del grupo al armar la exposición.
    Belinda estaba hablando muy bajo, por lo que le pidieron
    que repita. No pusieron a alguien más de colaborador
    para aceptar en la llamada. Mera explicó, pero se le
    notaba que estaba leyendo pausadamente. Loor estaba
    prestando atención más al celular que a su compañera Mera.
    Reinoso mencionó su proyecto, pero no dio una pequeña
    introducción para conocer el contexto. \\
    \hline

    D6 & Fuentes y créditos (10) & [0--4] &
    Todavia falta una de exponer un integrante\\
    \hline

\end{tabularx}

\textbf{Nota sugerida:}

\[
    \frac{
        10n_1+20n_2+15n_3+20n_4+25n_5+10n_6
    }{4}
    \times \frac{10}{100}
    = [x.x]
\]

\textbf{Fortaleza:} [Pendiente.]

\textbf{Mejora:} [Pendiente.]

\textbf{Hoja de campo:}
\path{hojas_campos/Grupo2 LMRT.jpeg}

% Anotación pendiente:
% Por ahora al grupo le falta exponer a un integrante
% por tiempo.

\clearpage

%============================================================
% DECLARACION DE IA
%============================================================

\section*{Declaración de uso de inteligencia artificial}

No se ha utilizado inteligencia artificial para la elaboración de este documento.

\end{document}